
Define monomial orderings give examples and state useful side definitions 

Monomial Orderings

A monomial ordering \(>\) on \(k[x_1, \ldots, x_n]\) is a relation \(>\) on \(\mathbb{Z}_{\geq 0}^n\), or equivalently, 
a relation on the set of monomials \(x^{\alpha}\), \(\alpha \in \mathbb{Z}_{\geq 0}^n\), satisfying:
1. \(>\) is a total (or linear) ordering on \(\mathbb{Z}_{\geq 0}^n\).
2. If \(\alpha > \beta\) and \(\gamma \in \mathbb{Z}_{\geq 0}^n\), then \(\alpha + \gamma > \beta + \gamma\).
3.  \item \(>\) is a well-ordering on \(\mathbb{Z}_{\geq 0}^n\). This means that every nonempty subset of \(\mathbb{Z}_{\geq 0}^n\) has a smallest element under \(>\). In other words, if \(A \subseteq \mathbb{Z}_{\geq 0}^n\) is nonempty, then there is \(\alpha \in A\) such that \(\beta > \alpha\) for every \(\beta \ne \alpha\) in \(A\).

An equivalent condition for well ordering is that \( \alpha \geq 0 \) for all \( \alpha \in \mathbb{Z}^n_{\geq 0} \).

Given a monomial ordering \(>\), we say that \(\alpha \geq \beta\) when either \(\alpha > \beta\) or \(\alpha = \beta\).

- Lexicographic (lex): Compare exponents left to right (\( \alpha > \beta \iff \alpha - \beta \) has its leftmost component positive.
- reverse lexicographic (rlex): Compare exponents right to left (\( \alpha > \beta \iff \alpha - \beta \) has its rightmost component negative.
- Graded lex (grlex): Compare total degree, then lex \( \alpha > \beta \iff |\alpha| > |\beta| \text{ or } |\alpha| = |\beta| \text{ and }\)
  \( \alpha >_{\text{lex}} \beta \).
- Graded reverse lex (grevlex): Compare total degree, then right to left.

Multidegree \( \text{multideg}(f) = \max \{ \alpha \in \mathbb{Z}^n_{\geq 0} | a_\alpha \neq 0 \} \) where \( f = \sum_\alpha a_{\alpha} x^\alpha \).
Leading monomial (LM): Largest monomial (w.r.t. order) with nonzero coefficient in \( f \).
Leading coefficient (LC): Coefficient of the LM in \( f \).
Leading term (LT): \( LC(f) \cdot LM(f) \).

The multidegree satisfies \( \text{multideg}(fg) = \text{multideg}(f) + \text{multideg}(g) \)
and \( \text{multideg}(f + g) \leq \max(\text{multideg}(f),\text{multideg}(g)) \), if \( f + g \neq 0 \).


Give the polynomial division algorithm with respect to some basis of an ideal


Input: \( f_1, \ldots, f_s, f \)
Output: \( q_1, \ldots, q_s, r \)

\[
q_1 := 0; \ldots; q_s := 0;\; r := 0
\]
\[
p := f
\]
WHILE \( p \ne 0 \) {
    \( i := 1 \)
    \( divisionoccurred := \mathrm{false} \)
    WHILE \( i \leq s \) AND \( divisionoccurred = \mathrm{false} \) {
        IF \( \mathrm{LT}(f_i) \) divides \( \mathrm{LT}(p) \) {
            \( q_i := q_i + \mathrm{LT}(p)/\mathrm{LT}(f_i) \)
            \( p := p - (\mathrm{LT}(p)/\mathrm{LT}(f_i)) f_i \)
            \( divisionoccurred := \mathrm{true} \)
        } ELSE {
            \( i := i + 1 \)
        }
    }
    IF \( divisionoccurred = \mathrm{false} \) {
        \( r := r + \mathrm{LT}(p) \)
        \( p := p - \mathrm{LT}(p) \)
    }
RETURN \( q_1, \ldots, q_s, r \)

The central difference compared to regular division is that the algorithm does not terminate if the intermediate 
dividend's \( p \) leading term is not divisible by any \( f_i \). Instead it moves the undivisble parts into the remainder 
term (which are not further divided by the \( f_i \) and just left as is).
E.g. \( x + y^2 + y : y^2 - 1 \) is turned into \( y^2 + y : y^2 - 1 \) and the remainder set to \( x \)
this becomes \( y + 1  \) which is no further divisible by \( y^2 - 1 \) and thus the final result reads:
\( f = 1\cdot (y^2 - 1) + (x + y + 1) \).



State some facts related to monomial ideals

An ideal \( I \subseteq k[x_1, \dots, x_n] \) is a monomial ideal if there is a subset
\( A \subseteq \mathbb{Z}^n_{\geq 0} \) (possibly infinite) such that \( I \) consists of all polynomials which are finite
sums of the form \( \sum_{\alpha \in A} h_{\alpha}x^\alpha \), where \( h_\alpha \in k[x_1, \dots, x_n] \). 
In this case, we write \( I = \langle x^\alpha | \alpha \in A \rangle \).

Lemma Let \( I = \langle x^\alpha | \alpha \in A \rangle  \) be a monomial ideal. 
Then \( x^\beta \in I \) if and only if \( x^\beta \) is divisible by \( x^\alpha \) for some \( \alpha \in A \).

Dickson's Lemma
Let \( I = \langle x^\alpha | \alpha \in A \rangle \) be a monomial ideal. 
Then \( I \) can be written in the form \( I = \langle x^{\alpha(1)}, \dots, x^{\alpha(s)} \rangle \), where
\( \alpha(1), \dots, \alpha(s) \in A \). In particular, I has a finite basis.
(Here \( \alpha(1), \alpha(s) \) are denoting different elements in \( A \).)

A monomial ideal \( I \subseteq k[x_1, \dots, x_n] \) has a basis \( x^{\alpha(1)}, \dots, x^{\alpha(s)} \) with
the property that \( x^{\alpha(i)} \) does not divide \( x^{\alpha(j)} \) for \( i \neq j \). 
Furthermore, this basis is unique and is called the minimal basis of I.


Define Groebner bases their reduced form and Buchberger’s Criterion

Groebner basis \( G = \{g_1,\ldots,g_t\} \) for an ideal \( I \): a finite generating set such that 
\( \langle LT(G) \rangle = \langle LT(I) \rangle \) where \( \langle LT(I) \rangle \) is the ideal of the leading terms of \( I \).
That is, for every \( f \in I \), \( LM(f) \) is divisible by \( LM(g_i) \) for some \( i \).

A Groebner basis is by definition a monomial ideal but is dependent on the monomial ordering chosen.
It can be shown that every Groebner basis is also a basis for an ideal \( I \).
Due to Hilbert's basis theorem a Groebner basis is always finite.
A Groebner basis is not unique but produces a unique remainder.

The main obstruction to reduced Groebner bases are Syzygy's that is if \( I = \langle LT(f_1), \dots, LT(f_n) \rangle \)
any collection of \( \{h_1, \dots, h_2\} \subset k[x_1, \dots, x_m] \) with
\[
    \sum_{i=1}^n h_i LT(f_i) = 0
\]
Those create issues since if we start our search for a basis by some set of leading terms of elements in \( I \)
their differences might not be again in this set (their leading terms cancel each other out).

For \( \text{multideg}(f) = \alpha, \text{multideg}(g) = \beta \), 
let \( \gamma = (\gamma_1, \dots, \gamma_m) \) with \( \gamma_i = \max(\alpha_i, \beta_i) \)
then \( x^\gamma := LCM(LM(f), LM(g)) \) is the least common multiple.
With that notion the insight of Buchberger was to add to the initial candidate set, remainders of elements of the form 
\( S(f, g) = \frac{x^\gamma}{LT(f)}\cdot f - \frac{x^\gamma}{LT(g)} \cdot g \).

Buchberger's criterion
Let \( I \) be a polynomial ideal. Then a basis
\( G = \{g_1, \dots, g_t \}  \) of \( I \) is a Gröbner basis of \( I \) if and only if for all pairs \( i \neq j \), the
remainder on division of \( S(g_i, g_j) \) by \( G \) (listed in some order) is zero.


Reduced Groebner basis: 
Unique for a fixed order; all leading coefficients are \( 1 \) and no monomial in any \( g_i \) 
is divisible by the leading monomial of any other \( g_j \).

A reduced Groebner basis is obtained by removing \( g_i \) from \( G \) dividing/reducing \( g_i \)
by the set \( G \setminus \{ g \} \) and adding the remainder to \( G \setminus \{ g_i \} \) until no progress is made.


Give applications for Groebner bases

Consider the equations

\[
x^2 + y^2 + z^2 = 1 \\
x^2 z^2 = y \\
x = z
\]
then a Groebner basis for the corresponding ideal of this variety is
\[
g_1 = x - z \\
g_2 = y - 2z^2 \\
g_3 = z^4 + \frac{1}{2} z^2 - \frac{1}{4}
\]
the last equation can be solved and inserting the solutions into the previous equations
gives all points (although an infinite amount) of the variety (note how much easier it is to determine the solutions in this shape).
One can consider this as a concrete instance of elimination theory.

Implicitization
Define the parameterized curve
\[
x = t^4
y = t^3
z = t^2
\]
Computing a Groebner basis of the corresponding ideal gives
\( G = \{ t^2 - z, ty - z^2, tz - y, x - z^2, y^2 - z^3 \} \)
then we know that the curve lies in the variety, defined by \( x - z^2 = 0, y^2 - z^3 = 0 \),
(which can be proven to be always the smallest variety containing it), which doesn't depend on the parameterization.
Again this is an instance of Elimination theory.

